\documentclass[12pt]{article}
\usepackage[T1]{fontenc}
\usepackage[utf8]{inputenc}
\usepackage{lmodern}
\usepackage{textcomp}
\usepackage{enumitem}
\usepackage{geometry}
\geometry{margin=1in}
\begin{document}
\begin{center}
Answer Key - Religious Studies - K-12 Level Assessment 2 \
\small (Answers are ordered exactly as in the question paper.)
\end{center}

\section*{Multiple Choice}
\begin{enumerate}[label=\arabic*.]
\item \textbf{Answer:} A
\\ \textit{Options:} A) Baptism and Eucharist; B) Baptism and Ordination; C) Eucharist and Marriage; D) Eucharist and Confession
\\ \textit{Note:} Baptism and Eucharist are recognized as the two essential sacraments in Christianity because they are universally practiced across most Christian denominations, symbolizing initiation into the faith and the commemorative act of Christ's Last Supper, respectively.
\item \textbf{Answer:} A
\\ \textit{Options:} A) Epic of Creation; B) Epic of Destruction; C) Epic of Egypt; D) Epic of Origins
\\ \textit{Note:} The Enuma Elish is commonly referred to as the "Epic of Creation" because it is a Babylonian creation myth that describes the origins of the world and the gods, detailing how the universe was formed and the establishment of order from chaos.
\item \textbf{Answer:} C
\\ \textit{Options:} A) Life stories; B) Warrior stories; C) Birth stories; D) Hero stories
\\ \textit{Note:} The Jataka tales are called "birth stories" because they recount the previous lives of the Buddha, illustrating moral lessons through his various incarnations.
\item \textbf{Answer:} A
\\ \textit{Options:} A) Mishnah; B) Midrash; C) Tanakh; D) Pentateuch
\\ \textit{Note:} The Oral Torah, which encompasses the interpretations and teachings that were initially transmitted orally, was eventually compiled and recorded in written form as the Mishnah, making it the correct answer.
\item \textbf{Answer:} A
\\ \textit{Options:} A) Akitu; B) Wag and Thoth; C) Bast; D) Nehebkau
\\ \textit{Note:} Akitu is the name of the ten-day New Year festival in ancient Babylon that celebrated the renewal of life and the culture of the Babylonians through various rituals and ceremonies honoring their gods.
\end{enumerate}

\section*{True / False}
\begin{enumerate}[label=\arabic*.]
\setcounter{enumi}{5}
\item \textbf{Answer:} True
\\ \textit{Options:} A) True; B) False
\\ \textit{Note:} The statement is true because zazen, which translates to "seated meditation," is a fundamental practice in Zen Buddhism that encourages practitioners to cultivate mindfulness and gain deeper insights into the nature of existence.
\item \textbf{Answer:} True
\\ \textit{Options:} A) True; B) False
\\ \textit{Note:} The Yijing, an ancient Chinese divination text, is accurately translated as The Classic of Changes due to its foundational role in exploring the principles of change and transformation in life.
\item \textbf{Answer:} False
\\ \textit{Options:} A) True; B) False
\\ \textit{Note:} The current Dalai Lama, Tenzin Gyatso, was born in 1935, not 1965.
\item \textbf{Answer:} True
\\ \textit{Options:} A) True; B) False
\\ \textit{Note:} Krishna and Rama are both revered in Hinduism as avatars of Vishnu, who is believed to take incarnations to restore cosmic order and protect dharma.
\item \textbf{Answer:} True
\\ \textit{Options:} A) True; B) False
\\ \textit{Note:} Francis of Assisi is recognized for receiving the stigmata, which are the marks resembling Christ's wounds, and for his well-documented sermons to birds, highlighting his deep connection with nature and animals.
\end{enumerate}

\section*{Long Form Response}
\begin{enumerate}[label=\arabic*.]
\setcounter{enumi}{10}
\item \textbf{Answer:} The Celestial Masters, a sect of Daoism founded by Zhang Daoling in the 2 nd century CE, emphasized the importance of rituals, community, and the worship of deities to achieve harmony with the Dao, or the natural way of the universe. Their beliefs included the idea of a cosmic order and the significance of moral purity, which they believed could lead to spiritual and physical well-being. The Celestial Masters became significant during the anti-Han rebellion, as their teachings offered hope and a sense of community to those disillusioned with the Han dynasty's rule. By promoting Daoist practices, they inspired social cohesion and resistance against corruption, reflecting how religion can serve both spiritual and political purposes in times of unrest.
\\ \textit{Note:} correct because it accurately describes the founding principles and practices of the Celestial Masters sect, including their emphasis on rituals and moral purity, and contextualizes their significance within the socio-political landscape of anti-Han rebellion in ancient China.
\item \textbf{Answer:} The construction of the first Buddhist temple in Japan in 596 CE marked a significant turning point in the country's cultural and religious landscape. This event introduced Buddhism, which would profoundly influence Japanese beliefs, art, and social structures. The temple served as a center for Buddhist practices and teachings, fostering the spread of new ideas and philosophies throughout Japan. As a result, many aspects of Japanese culture, such as architecture, literature, and rituals, began to reflect Buddhist principles. Ultimately, the establishment of this temple laid the foundation for Buddhism to become one of the major religions in Japan, shaping its spiritual identity for centuries to come.
\\ \textit{Note:} correct because it highlights the introduction of Buddhism as a transformative force in Japan, emphasizing its influence on culture, art, and social structures following the construction of the first temple, which became a hub for disseminating new ideas and practices.
\end{enumerate}

\vfill
\noindent\textit{End of Answer Key}
\end{document}